\documentclass[../general/general]{subfiles}

\begin{document}

\chapter{Лекция №8}
\noindent\textit{А.\,М.~Белемук \hfill 3 апреля 2021}
\vspace{1em}
\ifSubfilesClassLoaded{\tableofcontents\newpage}{}


% ===== СЕКЦИЯ 1: Переход из неравновесного состояния в равновесное =====
\section{Переход из неравновесного состояния в равновесное в изолированной системе}

Сегодня лекция 8. Мы продолжаем тему, начатую на прошлой лекции: рассматриваем переходы из неравновесного состояния в равновесное.

\subsection{Постановка задачи}

Рассмотрим изолированную систему, разбитую перегородками на две подсистемы. Каждая подсистема находится в состоянии локального равновесия --- это так называемое \emph{состояние частичного равновесия}. Общая энергия системы фиксирована: $E_1 + E_2 = E = \text{const}$.

Введём параметр отклонения от равновесия:
\begin{equation}
    x = E_1 - E_1^{\text{eq}}
\end{equation}
где $E_1^{\text{eq}}$ --- значение энергии первой подсистемы в состоянии равновесия. Параметр $x$ характеризует отклонение системы от равновесного состояния.

\subsection{Поведение энтропии}

Полная энтропия системы равна сумме энтропий подсистем (каждая находится в локальном равновесии):
\begin{equation}
    S = S_1(E_1) + S_2(E_2) \equiv S(E, x)
\end{equation}

По второму закону термодинамики, когда перегородку убирают, система самопроизвольно переходит в состояние полного равновесия $2$:
\begin{equation}
    \Delta S = S_2 - S_1^* \geq 0
\end{equation}

Следовательно, функция $S(x)$ монотонно возрастает при движении от неравновесного состояния к равновесному и достигает максимума при $x = 0$:
\begin{equation}
    S(x) \to \max \quad \text{при} \quad x \to 0
\end{equation}

В равновесии первая производная обращается в нуль:
\begin{equation}
    \left.\frac{\partial S}{\partial x}\right|_{x=0} = 0
\end{equation}

\subsection{Вероятность флуктуации в изолированной системе}

Рассмотрим обратный процесс: самопроизвольный уход системы из равновесия --- \textbf{флуктуацию}. Флуктуация --- это спонтанное отклонение системы от равновесия, характеризуемое параметром $x$.

Вероятность того, что система находится в неравновесном состоянии с параметром $x$, равна отношению числа микросостояний $\Omega(E, x)$ с данным $x$ к полному числу микросостояний $\Omega(E)$:
\begin{equation}
    W(x) = \frac{\Omega(E, x)}{\Omega(E)}
\end{equation}

Вводя энтропию неравновесного состояния $S(x) = \ln \Omega(E, x)$ по формуле Больцмана и учитывая, что $S^{\text{eq}} = \ln \Omega(E)$, получаем:
\begin{equation}
    \boxed{W(x) \propto e^{\Delta S}}
    \label{eq:W_isolated}
\end{equation}
где $\Delta S = S(x) - S^{\text{eq}} \leq 0$ --- изменение энтропии при флуктуации. Поскольку $\Delta S \leq 0$, вероятность флуктуации экспоненциально мала. Равновесие подавляет флуктуации.

Это \textbf{формула 1}: вероятность флуктуации в изолированной системе.

% ===== КОНЕЦ СЕКЦИИ 1 =====


% ===== СЕКЦИЯ 2: Тело в среде =====
\section{Тело в среде. Обобщение на открытые системы}

\subsection{Постановка задачи}

Рассмотрим тело, погружённое в среду (термостат). Вся система (тело + среда) изолирована. Среда находится в равновесии при любом процессе с телом и характеризуется фиксированными температурой $T_0$ и давлением $P_0$.

При переходе тела из неравновесного состояния $1^*$ в равновесное $2$ полная энтропия может только возрастать:
\begin{equation}
    \Delta S_{\text{total}} = \Delta S + \Delta S_{\text{среды}} \geq 0
\end{equation}
где $\Delta S$ (без индекса) --- изменение энтропии тела.

\subsection{Изменение энтропии среды}

Применим первый закон термодинамики к телу. Тепло $\Delta Q$, полученное телом в ходе перехода:
\begin{equation}
    \Delta Q = \Delta E + P_0\,\Delta V
\end{equation}
где $\Delta E$ --- изменение внутренней энергии тела, $P_0 \Delta V$ --- работа, совершённая телом над средой при постоянном давлении $P_0$.

Поскольку процесс в среде равновесный (среда всегда в термодинамическом равновесии), по второму закону:
\begin{equation}
    \Delta Q_{\text{среды}} = T_0\,\Delta S_{\text{среды}}
\end{equation}

Тепло, полученное средой, равно теплу, отданному телом: $\Delta Q_{\text{среды}} = -\Delta Q$. Поэтому:
\begin{equation}
    \Delta S_{\text{среды}} = \frac{-\Delta Q}{T_0} = \frac{-(\Delta E + P_0\,\Delta V)}{T_0}
\end{equation}

\subsection{Основное неравенство для тела в среде}

Подставляя в условие $\Delta S_{\text{total}} \geq 0$:
\begin{equation}
    \Delta S - \frac{\Delta E + P_0\,\Delta V}{T_0} \geq 0
\end{equation}

Умножим на $(-T_0) < 0$ (знак неравенства меняется):
\begin{equation}
    \boxed{\Delta E - T_0\,\Delta S + P_0\,\Delta V \leq 0}
    \label{eq:main_ineq}
\end{equation}

Это \textbf{ключевое неравенство}: при любом спонтанном переходе тела из неравновесного состояния в равновесное комбинация $\Delta E - T_0\Delta S + P_0\Delta V$ не возрастает.

\subsection{Следствие а: минимизация свободной энергии при фиксированных $T$, $V$}

Пусть объём тела фиксирован ($\Delta V = 0$) и тело находится в тепловом контакте со средой при $T_0$, причём температура тела не флуктуирует и равна $T_0$. Тогда $\Delta V = 0$, $T = T_0 = \text{const}$, и из~\eqref{eq:main_ineq}:
\begin{equation}
    \Delta E - T_0\,\Delta S \leq 0 \quad \Rightarrow \quad \Delta(E - TS) \leq 0 \quad \Rightarrow \quad \Delta F \leq 0
\end{equation}

\textbf{Вывод:} при фиксированных $T$ и $V$ свободная энергия Гельмгольца $F = E - TS$ может только убывать. В состоянии равновесия $F$ достигает \emph{минимума} по всем нефиксированным параметрам.

\subsection{Следствие б: минимизация потенциала Гиббса при фиксированных $T$, $P$}

Пусть теперь $T = T_0$ и $P = P_0$, то есть объём может меняться, но давление зафиксировано. Тогда из~\eqref{eq:main_ineq}:
\begin{equation}
    \Delta E - T_0\,\Delta S + P_0\,\Delta V \leq 0 \quad \Rightarrow \quad \Delta(E - TS + PV) \leq 0 \quad \Rightarrow \quad \Delta\Phi \leq 0
\end{equation}

\textbf{Вывод:} при фиксированных $T$ и $P$ потенциал Гиббса $\Phi = E - TS + PV$ может только убывать. В состоянии равновесия $\Phi$ достигает \emph{минимума}. Именно этот факт делает $\Phi$ главным объектом при описании фазовых переходов.

% ===== КОНЕЦ СЕКЦИИ 2 =====


% ===== СЕКЦИЯ 3: Термодинамические неравенства =====
\section{Термодинамические неравенства. Условия устойчивости}

Рассмотрим теперь обратную задачу: тело находится в равновесии, и мы рассматриваем \emph{спонтанные малые отклонения} от этого равновесия (флуктуации). При флуктуации неравенство~\eqref{eq:main_ineq} работает в обратную сторону: от равновесия $2$ к неравновесному состоянию $1^*$, поэтому:
\begin{equation}
    \delta(E - T_0\,S + P_0\,V) \geq 0
    \label{eq:stab_ineq}
\end{equation}

\subsection{Разложение по малым отклонениям}

Рассматриваем малые отклонения от равновесия. Считаем, что зависимость $E(S, V)$ в окрестности равновесия такая же, как в равновесии. Разложим $\Delta E$ по $\Delta S$ и $\Delta V$ до второго порядка:
\begin{equation}
    \Delta E = \left(\frac{\partial E}{\partial S}\right)_{\!V}\!\Delta S + \left(\frac{\partial E}{\partial V}\right)_{\!S}\!\Delta V + \frac{1}{2}\,d^2 E + \ldots
\end{equation}

Так как $(\partial E/\partial S)_V = T_0$ и $(\partial E/\partial V)_S = -P_0$, первый порядок в $\delta(E - T_0 S + P_0 V)$ сокращается. Остаётся второй порядок:
\begin{equation}
    \delta(E - T_0\,S + P_0\,V) = \frac{1}{2}\,d^2E \geq 0
\end{equation}

Раскрывая второй дифференциал:
\begin{equation}
    d^2 E = \frac{\partial^2 E}{\partial S^2}(\Delta S)^2 + 2\frac{\partial^2 E}{\partial S\,\partial V}\Delta S\,\Delta V + \frac{\partial^2 E}{\partial V^2}(\Delta V)^2 \geq 0
\end{equation}

Это квадратичная форма от двух переменных $(\Delta S,\, \Delta V)$ должна быть \textbf{положительно определённой}.

\subsection{Необходимые и достаточные условия положительной определённости}

По критерию Сильвестра квадратичная форма $a_{11}x_1^2 + 2a_{12}x_1 x_2 + a_{22}x_2^2$ положительно определена тогда и только тогда, когда:
\begin{equation}
    a_{11} > 0, \qquad \det\begin{pmatrix} a_{11} & a_{12} \\ a_{12} & a_{22} \end{pmatrix} > 0
\end{equation}

(условие $a_{22} > 0$ следует из второго условия автоматически).

\subsection{Первое термодинамическое неравенство: $C_V > 0$}

Из условия $a_{11} = \left(\partial^2 E/\partial S^2\right)_V > 0$:
\begin{equation}
    \left(\frac{\partial^2 E}{\partial S^2}\right)_{\!V} = \left(\frac{\partial T}{\partial S}\right)_{\!V} > 0
\end{equation}

Домножим числитель и знаменатель на $T$:
\begin{equation}
    \left(\frac{\partial T}{\partial S}\right)_{\!V} = \frac{T}{T\left(\partial S/\partial T\right)_V} = \frac{T}{C_V} > 0
\end{equation}

Поскольку $T > 0$, получаем первое термодинамическое неравенство:
\begin{equation}
    \boxed{C_V > 0}
    \label{eq:CV_pos}
\end{equation}

Теплоёмкость при постоянном объёме всегда положительна. Физический смысл: при подводе тепла к системе её температура должна возрастать.

\subsection{Из $a_{22} > 0$: адиабатическая сжимаемость}

Из $a_{22} = \left(\partial^2 E/\partial V^2\right)_S > 0$:
\begin{equation}
    \left(\frac{\partial^2 E}{\partial V^2}\right)_{\!S} = \left(\frac{\partial(-P)}{\partial V}\right)_{\!S} = -\left(\frac{\partial P}{\partial V}\right)_{\!S} > 0
\end{equation}

Вводя \emph{адиабатическую сжимаемость}:
\begin{equation}
    \kappa_S = -\frac{1}{V}\left(\frac{\partial V}{\partial P}\right)_{\!S}
\end{equation}
получаем:
\begin{equation}
    \boxed{\kappa_S > 0}
\end{equation}

\subsection{Из условия на детерминант: изотермическая сжимаемость и соотношение $C_P \geq C_V$}

Условие $\det(a_{ij}) > 0$ является наиболее содержательным. Раскрывая определитель матрицы вторых производных $E$ по переменным $(S, V)$ и переходя к переменным $(T, V)$, получаем:
\begin{equation}
    \det \begin{pmatrix} E_{SS} & E_{SV} \\ E_{SV} & E_{VV} \end{pmatrix} = \frac{T}{C_V}\cdot\frac{1}{V}\cdot\frac{1}{\kappa_T} \cdot\frac{T\,V}{C_V} \cdot \frac{C_V}{T\,\kappa_T}
\end{equation}

После алгебраических преобразований (перемножение определителей якобианов) условие $\det > 0$ сводится к двум дополнительным результатам.

\medskip
\noindent\textbf{Изотермическая сжимаемость:}
\begin{equation}
    \kappa_T = -\frac{1}{V}\left(\frac{\partial V}{\partial P}\right)_{\!T} > 0
    \label{eq:kT_pos}
\end{equation}

\textbf{Физический смысл:} при увеличении давления объём должен уменьшаться. Если $\kappa_T \leq 0$ --- система механически неустойчива.

\medskip
\noindent\textbf{Соотношение между сжимаемостями:}

Из тождества~\eqref{eq:kT_pos} и перестановки переменных в определителе получаем (аналогичным раскрытием с переменными $S, P$):
\begin{equation}
    \boxed{\kappa_T \geq \kappa_S}
    \label{eq:kT_ge_kS}
\end{equation}

Изотермическая сжимаемость не меньше адиабатической.

\medskip
\noindent\textbf{Соотношение теплоёмкостей:} используя термодинамическое тождество $C_P - C_V = TV\alpha^2/\kappa_T \geq 0$, получаем:
\begin{equation}
    \boxed{C_P \geq C_V > 0}
    \label{eq:CP_ge_CV}
\end{equation}

Таким образом, из условий положительной определённости квадратичной формы $d^2E$ вытекают четыре термодинамических неравенства:
\begin{equation}
    C_V > 0, \quad \kappa_S > 0, \quad \kappa_T > 0, \quad C_P \geq C_V, \quad \kappa_T \geq \kappa_S
\end{equation}

Это и есть \textbf{термодинамические условия устойчивости} равновесного состояния.

% ===== КОНЕЦ СЕКЦИИ 3 =====


% ===== СЕКЦИЯ 4: Флуктуации тела в среде =====
\section{Флуктуации тела в среде}

Теперь вернёмся к формуле~\eqref{eq:W_isolated} для вероятности флуктуации и применим её к телу в среде.

\subsection{Вероятность флуктуации тела в среде}

Рассматриваем тело + среда как изолированную систему. Флуктуация --- это состояние, при котором энергия, энтропия и объём тела отклонились от равновесных значений на $\Delta E$, $\Delta S$, $\Delta V$. Применяя общую формулу~\eqref{eq:W_isolated} и результат для $\Delta S_{\text{total}}$ из формулы~\eqref{eq:main_ineq}:
\begin{equation}
    \Delta S_{\text{total}} = \Delta S - \frac{\Delta E + P_0\,\Delta V}{T_0} = -\frac{1}{T_0}\left(\Delta E - T_0\,\Delta S + P_0\,\Delta V\right)
\end{equation}

Получаем вероятность флуктуации тела в среде:
\begin{equation}
    \boxed{W \propto \exp\!\left(-\frac{\Delta E - T_0\,\Delta S + P_0\,\Delta V}{T_0}\right)}
    \label{eq:W_medium}
\end{equation}

Комбинация $R = \Delta E - T_0\,\Delta S + P_0\,\Delta V$ называется \emph{работой Планка} (минимальная работа, необходимая для создания флуктуации). Формула~\eqref{eq:W_medium} --- это \textbf{формула Эйнштейна для флуктуаций}.

\subsection{Частные случаи}

При $\Delta V = 0$ (флуктуации при постоянном объёме):
\begin{equation}
    W \propto e^{-\Delta F/T_0}
\end{equation}

При $\Delta T = 0$ и $\Delta P = 0$ (флуктуации при постоянных $T$ и $P$):
\begin{equation}
    W \propto e^{-\Delta\Phi/T_0}
\end{equation}

\subsection{Квазиравновесные флуктуации}

Формула~\eqref{eq:W_medium} применима в режиме \textbf{квазиравновесных (квазитермодинамических) флуктуаций}. Введём два характерных времени:
\begin{itemize}
    \item $\tau_1$ --- характерное время обмена теплом между телом и средой (время термализации тела со средой),
    \item $\tau_2$ --- характерное время релаксации внутри тела (порядок времени межатомных столкновений).
\end{itemize}

Условие применимости:
\begin{equation}
    \tau_2 \ll \tau_1
\end{equation}

Это означает, что тело успевает прийти к локальному равновесию внутри себя (за время $\tau_2$) гораздо быстрее, чем обменяться энергией со средой (время $\tau_1$). Таким образом, тело всегда находится в локальном равновесии при каждом значении параметра флуктуации.

Наблюдение флуктуаций производится на масштабе времени $t_{\text{наблюд}} \gg \tau_1$, что позволяет усреднять по многим флуктуациям. На практике нас интересуют средние квадратичные флуктуации: $\langle(\Delta S)^2\rangle$, $\langle(\Delta V)^2\rangle$, $\langle(\Delta E)^2\rangle$.

На следующей лекции мы вычислим эти величины, используя формулу~\eqref{eq:W_medium}.

% ===== КОНЕЦ СЕКЦИИ 4 =====


\end{document}
